\documentclass{article}
\usepackage{amsmath} 
\usepackage{amsfonts}
\usepackage{amssymb}
\usepackage[parfill]{parskip}
\usepackage{esvect}
\usepackage[thinc]{esdiff}
\usepackage{mathtools}
\usepackage{physics}
\usepackage{amsthm}
\usepackage{xfrac}
\usepackage{enumitem}

\renewcommand*\diff{\mathop{}\!\mathrm{d}}

\begin{document}

\title{MATH 430: HW 6}
\author{Jacob Lockard}
\date{10 March 2026}

\maketitle

\section*{Exercise 5.60}

\begin{quote}
    \itshape
    Let $H$ be a subgroup of a group $G$. For $a, b \in G$,
    let $a \sim b$ if and only if $ab^{-1} \in H$.
    Show that $\sim$ is an equivalence relation on $G$.
\end{quote}

Let $a \in G$. Then $aa^{-1} = e \in H$, so $a \sim a$, and
$\sim$ is reflexive.
Now let $a, b \in G$ with $a \sim b$. Then,
\begin{align*}
    ba^{-1} = \bigl(ab^{-1}\bigr)^{-1} \in H \,,
\end{align*}
since $ab^{-1} \in H$ and $H$ is a subgroup.
So $b \sim a$, and $\sim$ is symmetric.
Finally, let $a, b, c \in G$ with $a \sim b$ and $b \sim c$. Then,
\begin{align*}
    ac^{-1} = aec^{-1} = a b^{-1} b c^{-1} \in H \,,
\end{align*}
since $ab^{-1} \in H$, $bc^{-1} \in H$, and $H$ is a subgroup.
So $a \sim c$, and $\sim$ is transitive.
\qedsymbol

\section*{Exercise 8.8}

\begin{quote}
    \itshape
    Determine whether the given map is a group homomorphism:

    ...

    Let $F$ be the additive group of all continuous functions mapping
    $\mathbb{R}$ to $\mathbb{R}$. Let $\phi : F \to \mathbb{R}$ be
    defined by $\phi(g) = \int_0^1 g(x) \diff x$.
\end{quote}

We'll show that $\phi$ is a homomorphism. Let $f, g \in F$.
We have:
\begin{align*}
    \phi(f + g) &= \int_0^1 (f+g)(x) \diff x \\
    &= \int_0^1 f(x) \diff x + \int_0^1 g(x) \diff x \\
    &= \phi(f) + \phi(g) \,,
\end{align*}
by linearity of the integration operator. \qedsymbol

\section*{Exercise 8.11}

\begin{quote}
    \itshape
    Compute the kernel for the given homomorphism $\phi$:

    ...

    $\phi : \mathbb{Z} \to \mathbb{Z}_8$ such that $\phi(1) = 6$
\end{quote}

We compute:
\begin{gather*}
    \phi(0) = 0 \,; \\
    \phi(1) = 6 \,; \\
    \phi(2) = 4 \,; \\
    \phi(3) = 2 \,; \\
    \phi(4) = 0 \,; \\
    ...
\end{gather*}

Apparently, $\ker(\phi) = \bigl\{ 4n \;\vert\; n \in \mathbb{N}_0 \bigr\}$.

\section*{Exercise 8.12}

\begin{quote}
    \itshape
    Compute the kernel for the given homomorphism $\phi$:

    ...

    $\phi : \mathbb{Z} \to \mathbb{Z}$ such that $\phi(1) = 12$
\end{quote}

Let $n \in \mathbb{Z}$. Then,
\begin{align*}
    \phi(n) = \phi(1 \cdot n) = n \phi(1) = 12n \,.
\end{align*}
So $\phi$ is clearly injective, and we conclude $\ker(\phi) = \{ e \}$.

\section*{Exercise 8.16}

\begin{quote}
    \itshape
    Compute the kernel for the given homomorphism $\phi$:

    ...

    Let $D$ be the additive group of all differentiable functions mapping
    $\mathbb{R}$ to $\mathbb{R}$ and $F$ the additive group of all functions
    from $\mathbb{R}$ to $\mathbb{R}$. $\phi : D \to f$ is given by
    $\phi(f) = f'$, the derivative of $f$.
\end{quote}

We know from calculus that if $f'$ is the identity function, then
$f$ has the form
\begin{align*}
    f(x) = x + c
\end{align*}
for some $c \in \mathbb{R}$. So $\ker(\phi)$ is the set of all such functions.

\section*{Exercise 10.2}

\begin{quote}
    \itshape
    Find all cosets of the subgroup $4\mathbb{Z}$ of $2\mathbb{Z}$.
\end{quote}

We have:
\begin{align*}
    0 + 4\mathbb{Z} &= 4 \mathbb{Z} \,; \\
    2 + 4\mathbb{Z} &= \bigl\{ 2 + 4n \;\vert\; n \in \mathbb{Z} \bigr\} \,; \\
    4 + 4\mathbb{Z} &= \bigl\{ 4 + 4n \;\vert\; n \in \mathbb{Z} \bigr\} = 4\mathbb{Z} \,; \\
    \dots
\end{align*}

\section*{Exercise 10.30}

\begin{quote}
    \itshape
    Let $H$ be a subgroup of a group $G$ such that $g^{-1}hg \in H$ for all
    $g \in G$ and all $h \in H$. Show that every left coset $gH$ is
    the same as the right coset $Hg$.
\end{quote}

Let $g \in G$. Let $gh \in gH$. Then we can write
\begin{align*}
    gh = (ghg^{-1}) g \,,
\end{align*}
where $ghg^{-1} = \bigl(g^{-1}\bigr)^{-1} h g^{-1} \in H$, since $g^{-1} \in G$.
So $gh \in Hg$. Similarly, if $hg \in Hg$, we can write
\begin{align*}
    hg = g(g^{-1}hg) \,,
\end{align*}
where $g^{-1} h g \in H$, so $hg \in gH$. We conclude $gH = Hg$. \qedsymbol

\section*{Exercise 10.33}

\begin{quote}
    \itshape
    Let $H$ be a subgroup of a group $G$ and let $a,b \in G$. Prove
    the statement or give a counterexample:

    ...

    If $Ha = Hb$, then $b \in Ha$.
\end{quote}

Assume $Ha = Hb$. Since $b = eb$, we have $b \in Hb = Ha$. \qedsymbol

\end{document}